\documentclass[10pt,a4paper]{amsart}
\usepackage[margin=19mm]{geometry}
\usepackage{amsmath,amssymb,mathtools}
\usepackage[scr=rsfs,cal=euler]{mathalfa}
\usepackage{xfrac}
\usepackage[unicode,hidelinks,bookmarks=false]{hyperref}
\newcommand{\ie}{i.e.\ }
\newcommand{\cf}{cf.\ }
\newcommand{\resp}{respectively\ }
\newcommand{\Subsection}{\subsection}
\providecommand{\nobreakdash}{\nobreak\hbox{-}}
\setlength{\emergencystretch}{2em}
\multlinegap0pt
% ---------------------------------------------------------------------------
% Editorial AMS wrapper prelude. The theorem environments and the mathematical macros below are
% transcribed from the paper's own preamble (cached-originals/source0.tex lines 1--97); the AMS amsart
% class, the named format packages of the pinned TeX Live runtime and this editorial formatting are not
% part of the author's text. No mathematics is changed.
% ---------------------------------------------------------------------------
\usepackage{cleveref}
\allowdisplaybreaks
\numberwithin{equation}{section}
\newtheorem{theorem}{Theorem}[section]
\newtheorem{proposition}[theorem]{Proposition}
\newtheorem{lemma}[theorem]{Lemma}
\newtheorem{corollary}[theorem]{Corollary}
\newtheorem{remark}[theorem]{Remark}
\newtheorem{example}[theorem]{Example}
\theoremstyle{definition}
\newtheorem{definition}[theorem]{Definition}
\newtheorem{Hypothesis}[theorem]{Hypothesis}
\DeclareMathOperator{\add}{add}
\DeclareMathOperator{\End}{End}
\DeclareMathOperator{\Aut}{Aut}
\DeclareMathOperator{\Hom}{Hom}
\DeclareMathOperator{\Ext}{Ext}
\DeclareMathOperator{\Coker}{Coker}
\DeclareMathOperator{\Ker}{Ker}
\DeclareMathOperator{\rad}{rad}
\DeclareMathOperator{\soc}{soc}
\DeclareMathOperator{\topm}{top}
\DeclareMathOperator{\Gr}{Gr}
\DeclareMathOperator{\ind}{ind}
\DeclareMathOperator{\coind}{coind}
\DeclareMathOperator{\proj}{proj}
\DeclareMathOperator{\inj}{inj}
\DeclareMathOperator{\modu}{mod}
\newcommand{\C}{\mathcal C}
\newcommand{\A}{\mathcal A}
\newcommand{\T}{\mathcal T}
\newcommand{\Pcal}{\mathcal P}
\newcommand{\E}{\mathbb E}
\newcommand{\F}{\mathsf F}
\newcommand{\X}{\mathsf X}
\newcommand{\kk}{\mathbb C}
\newcommand{\Ksp}{K_0^{\mathrm{sp}}}
\newcommand{\underlineC}{\underline{\mathcal C}}
\newcommand{\underlineHom}{\underline{\Hom}}
\newcommand{\xto}[1]{\xrightarrow{#1}}
\newcommand{\etri}{\dashrightarrow}
\newcommand{\bx}{\mathbf x}
\newcommand{\by}{\mathbf y}
\newcommand{\be}{\mathbf e}
\newcommand{\bdim}{\underline{\dim}}
\newcommand{\ol}[1]{\overline{#1}}

\begin{document}
\title{A refined multiplication formula in 2-Calabi-Yau Frobenius extriangulated categories}
\author{Ming Ding, Fan Xu and Panyue Zhou}
\date{Source: arXiv:2609.20049v1; distributed under CC BY 4.0}
\maketitle
\noindent\emph{Source, version and licence.} \emph{A refined multiplication formula in 2-Calabi-Yau Frobenius extriangulated categories}, by Ming Ding, Fan Xu and Panyue Zhou. The source of this editable unit is the e-print of \textbf{version v1} of arXiv:2609.20049, https://arxiv.org/abs/2609.20049v1, whose material is distributed under the Creative Commons Attribution 4.0 International licence, http://creativecommons.org/licenses/by/4.0/, as recorded in the arXiv licence metadata for this paper and shown on that abstract page. The whole of the body below is version v1, the newest version of the paper. Full rights notice: licenses/arxiv-2609.20049-CC-BY-4.0.txt.\par\smallskip
\noindent\textit{Editorial scope.} Counted unit: original Theorem 3.9 of the article, the statement carried by the source environment \texttt{theorem} with source label \texttt{thm:main}, cached-originals/source0.tex lines 662--675, printed as \textbf{Theorem 3.9} exactly as in the article: the multiplication formula that expresses the product in the Ext-algebra of a 2-Calabi-Yau Frobenius extriangulated category in terms of the splitting of the extensions into the two parts of the associated lattice, as stated in the paper. It is delivered together with the authors' own complete proof as the single primary proof body, source0.tex lines 677--719 (43 lines): the \texttt{proof} environment that belongs to the statement and establishes it in full. No proof marker is added and no mathematics is written, repaired or re-derived; the authors' notation and the printed slips of the source are kept literal. Necessary complete context is retained uncounted, in source order: lines 319--323, the source \texttt{equation} with source label \texttt{eq:character} of the article that the retained passages refer to; lines 434--436, the source \texttt{Hypothesis} with source label \texttt{hyp:constructible-cones} of the article that the retained passages refer to; lines 489--496, the source \texttt{proposition} with source label \texttt{prop:stratification} of the article that the retained passages refer to; lines 525--529, the source \texttt{equation} with source label \texttt{eq:lhs-expansion} of the article that the retained passages refer to; lines 553--559, the source \texttt{lemma} with source label \texttt{lem:forward-fibres} of the article that the retained passages refer to; lines 596--602, the source \texttt{lemma} with source label \texttt{lem:complement} of the article that the retained passages refer to; lines 625--628, the source \texttt{equation} with source label \texttt{eq:forward-exponent} of the article that the retained passages refer to; lines 630--633, the source \texttt{equation} with source label \texttt{eq:reverse-exponent} of the article that the retained passages refer to. The article's own numbering is reproduced by explicit counter settings: the theorem-like counter of the article is set before each retained passage, so the counted statement prints with the number the article gives it. Every \texttt{\textbackslash ref} and \texttt{\textbackslash eqref} inside every retained passage resolves inside this delivered file, and the 1 citation(s) that occur in the retained passages, \texttt{Palu2012}, are delivered as the authors' own bibliography entries transcribed verbatim from the article's own bibliography, keeping the authors' own printed numbers. The source's own class is \texttt{article}; the e-print ships no class or style file, so no publisher class is shipped, used or redistributed, and the wrapper is the AMS \texttt{amsart} class with the article's own macros and theorem declarations transcribed from its preamble. The retained passages contain no figure environment and no graphics-inclusion command, so no artwork is redistributed. Every declared body of this unit is an exact, unmodified span of source0.tex: no body is a bounded passage comparison, \texttt{omit} was not used anywhere, and consequently no fidelity receipt is asserted in delivery-evidence.json. The complete concatenated original source is bundled as sources/210832/source0.tex. Omitted from this editable unit and therefore absent from it are source0.tex lines 1--318; 324--433; 437--488; 497--524; 530--552; 560--595; 603--624; 629--629; 634--661; 676--676; 720--1453; 1458--1511. The AMS \texttt{amsart} wrapper, the named format packages of the pinned TeX Live runtime and this editorial source, licence and scope paragraph are editorial formatting only.\par\smallskip
\setcounter{section}{2}
\setcounter{equation}{2}
\setcounter{theorem}{1}
\begin{equation}\label{eq:character}
 \X^T_M
 =\bx^{\ind_T(M)}
   \sum_e\chi\bigl(\Gr_e(\F M)\bigr)\bx^{-\Phi(e)}.
\end{equation}

\setcounter{section}{3}
\setcounter{equation}{3}
\setcounter{theorem}{0}
\begin{Hypothesis}\label{hyp:constructible-cones}
The stable category $\A=\underline{\C}$ has constructible cones with respect to $T'$ in the sense of Palu \cite[Section~1.3]{Palu2012}. Equivalently, for every pair $L,M\in\C$, the cylinders over $\pi L\to\Sigma\pi M$ and $\pi M\to\Sigma\pi L$ are constructible with respect to $T'$ in the preceding sense.
\end{Hypothesis}

\setcounter{section}{3}
\setcounter{equation}{4}
\setcounter{theorem}{3}
\begin{proposition}\label{prop:stratification}
Under \cref{hyp:constructible-cones}, the decomposition
\[
 \E(L,M)=\bigsqcup_{Y\in\mathcal Y_{L,M}}
 \E(L,M)_{\langle Y\rangle}
\]
is a finite constructible stratification.
\end{proposition}

\setcounter{section}{3}
\setcounter{equation}{4}
\setcounter{theorem}{5}
\begin{equation}\label{eq:lhs-expansion}
 \chi(\mathbb PV)\X^T_L\X^T_M
 =\bx^{\ind_T(L\oplus M)}
   \sum_{e,f}\chi(\mathcal B_{e,f})\bx^{-\Phi(e+f)}.
\end{equation}

\setcounter{section}{3}
\setcounter{equation}{5}
\setcounter{theorem}{5}
\begin{lemma}\label{lem:forward-fibres}
Every nonempty fibre of $\Psi_{e,f}$ is an affine space.  Consequently,
\[
 \chi\bigl(\mathcal L^V_1(e,f)\bigr)
 =\sum_g\chi\bigl(\mathcal W^V_{L,M}(e,f,g)\bigr).
\]
\end{lemma}

\setcounter{section}{3}
\setcounter{equation}{5}
\setcounter{theorem}{6}
\begin{lemma}\label{lem:complement}
For every $e,f$, we have
\[
 \chi\bigl(\mathcal L^V_2(e,f)\bigr)
 =\sum_g\chi\bigl(\mathcal W^{\mathcal R}_{M,L}(f,e,g)\bigr).
\]
\end{lemma}

\setcounter{section}{3}
\setcounter{equation}{5}
\setcounter{theorem}{8}
\begin{equation}\label{eq:forward-exponent}
 \ind_T(L\oplus M)-\Phi(e+f)
 =\ind_T(Y_\delta)-\Phi(g).
\end{equation}

\setcounter{section}{3}
\setcounter{equation}{6}
\setcounter{theorem}{8}
\begin{equation}\label{eq:reverse-exponent}
 \ind_T(L\oplus M)-\Phi(e+f)
 =\ind_T(Y'_\eta)-\Phi(g).
\end{equation}

\setcounter{section}{3}
\setcounter{equation}{7}
\setcounter{theorem}{8}
\begin{theorem}\label{thm:main}
Assume \cref{hyp:constructible-cones}.  Let $L,M\in\C$ and let
$0\neq V\subseteq\E(L,M)$ be a linear subspace.  Then
\begin{equation}\label{eq:main}
\begin{aligned}
 \chi(\mathbb PV)\X^T_L\X^T_M
 ={}&\sum_{Y\in\mathcal Y_{L,M}}
 \chi\bigl(\mathbb P(V\cap\E(L,M)_{\langle Y\rangle})\bigr)\X^T_Y\\
 &+\sum_{Y\in\mathcal Y_{M,L}}
 \chi\bigl(\mathcal R_V\cap
   \mathbb P\E(M,L)_{\langle Y\rangle}\bigr)\X^T_Y.
\end{aligned}
\end{equation}
\end{theorem}

\setcounter{section}{3}
\setcounter{equation}{8}
\setcounter{theorem}{9}
\begin{proof}
Start with the expansion \eqref{eq:lhs-expansion}.  Since
\[
 \mathcal B_{e,f}=\mathcal L^V_1(e,f)\sqcup\mathcal L^V_2(e,f),
\]
additivity of Euler characteristic gives
\begin{align*}
 \chi(\mathbb PV)\X^T_L\X^T_M
 ={}&\bx^{\ind_T(L\oplus M)}
 \sum_{e,f}\chi\bigl(\mathcal L^V_1(e,f)\bigr)\bx^{-\Phi(e+f)}\\
 &+\bx^{\ind_T(L\oplus M)}
 \sum_{e,f}\chi\bigl(\mathcal L^V_2(e,f)\bigr)\bx^{-\Phi(e+f)}.
\end{align*}

By \cref{lem:forward-fibres}, the first term is
\[
 \sum_{e,f,g}\chi\bigl(\mathcal W^V_{L,M}(e,f,g)\bigr)
 \bx^{\ind_T(L\oplus M)-\Phi(e+f)}.
\]
Decompose each incidence variety according to the strata of middle terms of \cref{prop:stratification}.  On the stratum indexed by $Y$, \eqref{eq:forward-exponent} replaces the monomial by
\[
 \bx^{\ind_T(Y)-\Phi(g)}.
\]
For fixed $[\delta]$ and $U\subseteq\F Y_\delta$, the classes $e=[b_\delta(U)]$ and $f=[a_\delta^{-1}(U)]$ are uniquely determined.  Hence summing over $e$ and $f$ simply forgets these two labels and recovers the quiver Grassmannian of the middle module.  The first term is therefore
\[
 \sum_{Y\in\mathcal Y_{L,M}}
 \chi\bigl(\mathbb P(V\cap\E(L,M)_{\langle Y\rangle})\bigr)
 \bx^{\ind_T(Y)}
 \sum_g\chi\bigl(\Gr_g(\F Y)\bigr)\bx^{-\Phi(g)}.
\]
By \eqref{eq:character}, this is the first sum on the right side of \eqref{eq:main}.

For the second term, \cref{lem:complement} replaces
$\mathcal L^V_2(e,f)$ by the incidence varieties of reverse extensions in $\mathcal R_V$.  The reverse exponent identity \eqref{eq:reverse-exponent} then gives
\[
 \sum_{Y\in\mathcal Y_{M,L}}
 \chi\bigl(\mathcal R_V\cap
   \mathbb P\E(M,L)_{\langle Y\rangle}\bigr)
 \bx^{\ind_T(Y)}
 \sum_g\chi\bigl(\Gr_g(\F Y)\bigr)\bx^{-\Phi(g)}.
\]
This is the second sum in \eqref{eq:main}.
\end{proof}

\begin{thebibliography}{99}
\bibitem[17]{Palu2012} Y.~Palu, Cluster characters II: A multiplication formula, \emph{Proc. Lond. Math. Soc.} \textbf{104} (2012).
\end{thebibliography}
\end{document}
